% SGA 3，第 VII_B 篇，英文重建稿——第 1 轮。
% 范围：权威底本分篇第 96--97 页，第 5.5 节全文，自开头的定义至注 5.5.2(b)。
% 合并阅读本页码： 609--610.
% 再版印刷页码： 598--599.
% 本范围内编者注： 169.
% 权威底本： Polo--Gille Exp7B-13oct24.pdf,
% SHA-256 D303690E0195654BBE7C3DAA5D89EA23E24CD709A3575217234AC19AE8836694.
% 曾参考 Jacob Reinhold 的 CC BY 4.0 Markdown 作为比对谱系。
% 本组件没有图表；正合列以原生 LaTeX 重建。

\subsection*{5.5. 补充}
\phantomsection
\label{sga3:VIIB:5.5}
\addcontentsline{toc}{subsection}{5.5. 补充}
\footnote{\label{sga3:VIIB:nde:169}编者增加了本小节，以给出第 III 和 VIA 讲中提到的 \SGARefLink{sga3:vii:corollary:5-4}{5.4} 的一些推论。}

回顾下列定义。若 Noether 局部环 \(A\) 的完备化 \(\widehat A\) 是一个完备正则 Noether 局部环 \(B\) 对由 \(B\) 中的一个正则序列生成的理想 \(I\) 的商，则称此环为\emph{完全交}\textup{（参见 EGA IV\(_4\)，19.3.1）}。

另一方面，设
\[
  \tau:Y\hookrightarrow X
\]
为概形的闭浸入。若 \(y\in Y\)，则当
\[
  \mathcal O_{X,y}\longrightarrow\mathcal O_{Y,y}
\]
的核由正则序列生成时，称 \(\tau\) \emph{在点 \(y\) 处为正则浸入}。若再有 \(X\) 局部 Noether，且 \(\tau\) 在每一点处都是正则浸入，就称 \(\tau\) 为\emph{正则浸入}\textup{（参见\textit{同上}，命题~16.9.10 和定义~16.9.2）}。

\medskip
\noindent\textbf{推论 5.5.1.}
\phantomsection
\label{sga3:VIIB:5.5.1}
\textit{若 \(G\) 为域 \(k\) 上的代数群，则局部环 \(\mathcal O_{G,e}\) 为完全交。}

\medskip
\noindent\textit{证明。}
由 EGA IV\(_4\)，19.3.4，可假设 \(k\) 代数闭。若 \(\operatorname{char}(k)=0\)，则已知 \(G\) 平滑\textup{（参见 \SGARefLink{sga3:VIIA:3.3.1}{3.3.1} 或 VIB，\SGARefLink{cor:sga3-VIB-1.6.1}{1.6.1}）}；因此，由 EGA IV\(_4\)，17.5.3 \((\mathrm d^{\prime\prime})\)，\(\widehat{\mathcal O}_{G,e}\) 为 \(k\) 上的形式幂级数代数。若 \(\operatorname{char}(k)=p>0\)，则将 \SGARefLink{sga3:vii:corollary:5-4}{5.4} 用于 \(H=\{e\}\)，即知 \(\widehat{\mathcal O}_{G,e}\) 为完全交。

\medskip
\noindent\textbf{注 5.5.2.}
\phantomsection
\label{sga3:VIIB:5.5.2}
\textit{设 \(k\) 为域，\(G\) 为平滑代数 \(k\)-群，\(H\) 为 \(G\) 的闭子群。}

\smallskip
\noindent\textit{(a)}
在 \SGARefLink{sga3:iii:proposition:4-15}{第 III 篇，4.15} 中，我们已见浸入
\[
  H\hookrightarrow G
\]
是正则的。这也可以如下由 \SGARefLink{sga3:vii:corollary:5-4}{5.4} 推出。如\textit{同上}，可假设 \(k\) 代数闭，且只须证明
\[
  \mathcal O_{G,e}\longrightarrow\mathcal O_{H,e}
\]
的核 \(I\) 由正则序列生成。置
\[
  A=\mathcal O_{G,e},
  \qquad
  \widehat G=\operatorname{Spf}(\widehat A),
  \qquad
  \widehat H
  =
  \operatorname{Spf}(\widehat{\mathcal O}_{H,e}).
\]
由于 \(A\) 为 Noether 环，有正合序列
\[
  0
  \longrightarrow I\otimes_A\widehat A
  \longrightarrow \widehat A
  \xrightarrow{\ \pi\ }
  \mathscr A(\widehat H)
  \longrightarrow 0,
\]
且 \(\widehat A\) 在 \(A\) 上忠实平坦。因此，由 EGA IV\(_4\)，16.9.10(ii) 和 19.1.5(ii)，只须证明 \(\pi\) 的核
\[
  \widehat I=I\otimes_A\widehat A
\]
由 \(\widehat A\) 中的一个正则序列生成。

由于 \(G\) 平滑，\(\widehat A\) 约化。因此，由 \SGARefLink{sga3:vii:corollary:5-4}{5.4}，子代数
\[
  B=\mathscr A(\widehat G/\widehat H)
\]
同构于形式幂级数代数
\[
  k[[x_1,\ldots,x_n]].
\]
从而，\(\widehat G/\widehat H\) 的单位截面在 \(B\) 中由正则序列 \((x_1,\ldots,x_n)\) 定义。由于 \(\widehat A\) 为 Noether 环，\(\widehat A\) 中由 \(x_1,\ldots,x_n\) 生成的理想 \(J\) 是闭的，因而由推论~1.4 等于 \(\widehat I\)。此外，\(\widehat A\) 拓扑平坦，从而在 \(B\) 上平坦\textup{（参见 \SGARefLink{sga3:VIIB:0.3.8}{0.3.8}）}。因此，由 EGA IV\(_4\)，19.1.5(ii)，\((x_1,\ldots,x_n)\) 为 \(\widehat A\) 中的正则序列。

\smallskip
\noindent\textit{(b)}
由 \SGARefLink{sga3:VIIB:5.2:theorem:item:ii}{5.2(ii)} 还可推出下述更精确的结果。假设 \(k\) 为完美域。将 \SGARefLink{sga3:VIIB:5.2:theorem:item:ii}{5.2(ii)} 用于代数
\[
  C=\mathscr A(H),
\]
可知存在 \(\omega_H\) 的一组基
\[
  (y_1,\ldots,y_{r+s})
\]
以及整数
\[
  1\leq n_1\leq\cdots\leq n_r
\]
使 \(\mathscr A(H)\) 同构于
\[
  k[[y_1,\ldots,y_{r+s}]]
\]
对由
\[
  y_i^{p^{n_i}},
  \qquad i=1,\ldots,r.
\]
生成的理想的商。将各 \(y_i\) 提升为 \(\omega_G\) 的元素 \(x_i\)，并将 \((x_1,\ldots,x_{r+s})\) 扩充为 \(\omega_G\) 的一组基 \((x_1,\ldots,x_n)\)。由于 \(\mathscr A(G)\) 约化，由 \SGARefLink{sga3:VIIB:5.2:theorem:item:iii}{5.2(iii)}，态射
\[
  k[[x_1,\ldots,x_n]]
  \longrightarrow
  \mathscr A(G)
\]
为同构。

因此，\(G\) 上存在一个“坐标系”\((x_1,\ldots,x_n)\)，即一个同构
\[
  \mathscr A(G)
  \simeq
  k[[x_1,\ldots,x_n]],
\]
使 \(H\) 由方程
\[
  x_i^{p^{n_i}}=0
  \quad (i=1,\ldots,r),
  \qquad
  x_i=0
  \quad (i>r+s).
\]
定义。这就是“Dieudonn\'e 定理”；参看 \textup{[Di55]}，\(\S\)19，定理~6，以及 \textup{[Di73]}，II，\(\S\)\SGARefLink{sga3:VIIB:3.2:prop}{3.2}，命题~3 及其前面的内容。

% 下一游标：权威底本分篇第 97 页的 VII_B 末尾参考文献，合并阅读本第 610 页，再版印刷第 599 页。
